\documentclass[../../main.tex]{subfiles}% 注意这里的文件路径不能用 ./main.tex ,否则用latexmk编译子文件会报错
\graphicspath{{\subfix{./image/}}{\subfix{../../image/}}} % 指定图片引用路径,后续可以直接使用图片文件名
% 如果图片存放在上级目录,则使用 ../../image/ 作为备用路径

% 例如:
% \begin{figure}[H]
% \centering
% \includegraphics[width=0.3\textwidth]{图.png}
% \caption{}
% \label{figure:图}
% \end{figure}
% 注意:上述\label{}一定要放在\caption{}之后,否则引用图片序号会只会显示??.

\begin{document}

\section{测度}\label{section:ch01-测度}

\subsection{一般性质}\label{subsection:ch01-测度一般性质}

\begin{definition}\label{definition:ch01-集合函数性质}
设 $\varnothing \in\mathcal{F}\subseteq\mathcal{P}(X)$。集合函数 $\varphi:\mathcal{F}\to[0,\infty]$ 称为：
\begin{enumerate}[(i)]
\item 单调的，若对任意 $E,F\in\mathcal{F}$、$E\subseteq F$ 有 $\varphi(E)\leqslant\varphi(F)$；
\item 可加的，若对任意有限不交族 $\{E_1,\cdots ,E_n\}\subseteq\mathcal{F}$ 且 $\bigcup_{k=1}^n E_k\in\mathcal{F}$ 有
\[
\varphi\left(\bigcup_{k=1}^n E_k\right)=\sum_{k=1}^n\varphi(E_k);
\]
\item $\sigma$--次可加的，若对任意序列 $\{E_k\}\subseteq\mathcal{F}$ 且 $\bigcup_{k=1}^\infty E_k\in\mathcal{F}$ 有
\[
\varphi\left(\bigcup_{k=1}^\infty E_k\right)\leqslant\sum_{k=1}^\infty\varphi(E_k);
\]
\item $\sigma$--可加的，若对任意不交序列 $\{E_k\}\subseteq\mathcal{F}$ 且 $\bigcup_{k=1}^\infty E_k\in\mathcal{F}$ 有
\[
\varphi\left(\bigcup_{k=1}^\infty E_k\right)=\sum_{k=1}^\infty\varphi(E_k).
\]
\end{enumerate}
\end{definition}

\begin{definition}\label{definition:ch01-测度}
设 $(X,\mathcal{A})$ 为一可测空间。集合函数 $\mu:\mathcal{A}\to[0,\infty]$ 称为 $\mathcal{A}$ 上（或 $X$ 上）的一个（正）测度，若 $\mu(\varnothing )=0$ 且 $\mu$ 为 $\sigma$--可加的。三元组 $(X,\mathcal{A},\mu)$ 称为一个测度空间。

测度 $\mu$ 称为有限的，若 $\mu(X)<\infty$；它称为 $\sigma$--有限的，若存在序列 $\{E_k\}\subseteq\mathcal{A}$ 使得 $X=\bigcup_{k=1}^\infty E_k$ 且对所有 $k\in\mathbb{N}$ 有 $\mu(E_k)<\infty$。测度空间 $(X,\mathcal{A},\mu)$ 称为有限的（相应地，$\sigma$--有限的），若 $\mu$ 为有限的（相应地，$\sigma$--有限的）。

若 $\mu(X)=1$，则测度 $\mu$ 称为概率测度，空间 $(X,\mathcal{A},\mu)$ 称为概率空间。$X$ 上概率测度的集合记为 $\mathcal{P}(X)$。

若 $\mathcal{F}\subseteq\mathcal{P}(X)$ 为任意子集族，则上述定义可推广如下。映射 $\mu:\mathcal{F}\to[0,\infty]$ 是 $\mathcal{F}$ 上的一个测度，若当 $\varnothing \in\mathcal{F}$ 时 $\mu(\varnothing )=0$，且对任意不交序列 $\{E_k\}\subseteq\mathcal{F}$ 满足 $\bigcup_{k=1}^\infty E_k\in\mathcal{F}$，有
\[
\mu\left(\bigcup_{k=1}^\infty E_k\right)=\sum_{k=1}^\infty\mu(E_k).
\]
$\mathcal{F}$ 上的测度 $\mu$ 称为 $\sigma$--有限的，若存在序列 $\{E_k\}\subseteq\mathcal{F}$ 使得 $X=\bigcup_{k=1}^\infty E_k$ 且对每个 $k\in\mathbb{N}$ 有 $\mu(E_k)<\infty$。
\end{definition}

\begin{remark}\label{remark:ch01-测度子空间}
设 $(X,\mathcal{A},\mu)$ 为一测度空间，$F\in\mathcal{A}$。$\mu$ 在迹 $\sigma$--代数 $\mathcal{A}\cap F$ 上的限制 $\mu|_{\mathcal{A}\cap F}$ 是由 $\mu$ 在 $\mathcal{A}\cap F$ 上诱导的测度。测度空间 $(F,\mathcal{A}\cap F,\mu|_{\mathcal{A}\cap F})$ 称为 $(X,\mathcal{A},\mu)$ 的一个测度子空间。
\end{remark}

以下测度的性质容易证明。

\begin{proposition}\label{proposition:ch01-测度性质}
设 $(X,\mathcal{A},\mu)$ 为一测度空间。则：
\begin{enumerate}[(i)]
\item 对任意有限不交族 $\{E_1,\cdots ,E_n\}\subseteq\mathcal{A}$，有
\[
\mu\left(\bigcup_{k=1}^n E_k\right)=\sum_{k=1}^n\mu(E_k);
\]
\item 对任意 $E,F\in\mathcal{A}$，$E\subseteq F\Rightarrow\mu(E)\leqslant\mu(F)$；
\item 对任意序列 $\{E_k\}\subseteq\mathcal{A}$，有
\[
\mu\left(\bigcup_{k=1}^\infty E_k\right)\leqslant\sum_{k=1}^\infty\mu(E_k);
\]
\item 对任意非减序列 $\{E_k\}\subseteq\mathcal{A}$，有
\[
\mu\left(\bigcup_{k=1}^\infty E_k\right)=\lim_{k\to\infty}\mu(E_k);
\]
\item 对任意非增序列 $\{E_k\}\subseteq\mathcal{A}$ 且 $\mu(E_1)<\infty$，有
\[
\mu\left(\bigcap_{k=1}^\infty E_k\right)=\lim_{k\to\infty}\mu(E_k).
\]
\end{enumerate}
\end{proposition}

\begin{example}\label{example:ch01-计数与Dirac}
(i) 设 $X$ 为一集合。映射 $\mu^\#:\mathcal{P}(X)\to[0,\infty]$ 定义为
\[
\mu^\#(E)\triangleq 
\begin{cases}
E\text{ 的元素个数} & \text{若 }E\text{ 为有限集},\\
\infty & \text{否则}
\end{cases}
\]
是一个测度，称为计数测度。显然，若 $X$ 可数，则 $\mu^\#$ 是 $\sigma$--有限的。

(ii) 设 $X\neq\varnothing $，$x\in X$。映射 $\delta_x:\mathcal{P}(X)\to[0,\infty)$ 定义为
\[
\delta_x(E)\triangleq 
\begin{cases}
1 & \text{若 }x\in E,\\
0 & \text{否则}
\end{cases}
\]
是一个测度，称为集中于 $\{x\}$ 的 Dirac 测度。
\end{example}

\subsection{Borel 测度与 Radon 测度}\label{subsection:ch01-Borel测度}

设 $(X,\mathcal{T})$ 为一 Hausdorff 空间，$\mathcal{B}=\mathcal{B}(X)$ 为由拓扑 $\mathcal{T}$ 生成的 Borel $\sigma$--代数。设 $\mathcal{K}\subseteq\mathcal{P}(X)$ 表示紧子集族。

\begin{definition}\label{definition:ch01-Borel测度}
\label{remark:ch01-局部有限}
设 $\mathcal{A}\supseteq\mathcal{B}$ 为一 $\sigma$--代数，$\mu:\mathcal{A}\to[0,\infty]$ 为一测度。
\begin{enumerate}[(a)]
\item $\mu$ 称为局部有限的，若对任意 $x\in X$，存在 $x$ 的一个邻域 $U$ 使得 $\mu(U)<\infty$。局部有限的测度 $\mu:\mathcal{A}\to[0,\infty]$ 称为 $X$ 上的一个 Borel 测度。
\item 集合 $E\in\mathcal{A}$ 称为 $\mu$--外正则的（或外正则的），若
\[
\mu(E)=\inf\{\mu(A)\mid A\supseteq E,\ A\in\mathcal{T}\}.
\]
测度 $\mu$ 称为外正则的，若每个 $E\in\mathcal{A}$ 都是 $\mu$--外正则的。
\item 集合 $E\in\mathcal{A}$ 称为 $\mu$--内正则的（或内正则的），若
\[
\mu(E)=\sup\{\mu(K)\mid K\subseteq E,\ K\in\mathcal{K}\}.
\]
测度 $\mu$ 称为内正则的，若每个 $E\in\mathcal{A}$ 都是 $\mu$--内正则的。
\item 集合 $E\in\mathcal{A}$ 称为 $\mu$--正则的（或正则的），若它既是 $\mu$--外正则的又是 $\mu$--内正则的。测度 $\mu$ 称为正则的，若每个 $E\in\mathcal{A}$ 都是 $\mu$--正则的。
\item 内正则的 Borel 测度称为 Radon 测度。
\item Borel 测度称为适度的，若存在 $\{A_n\}\subseteq\mathcal{T}$ 使得 $X=\bigcup_{n=1}^\infty A_n$ 且对所有 $n\in\mathbb{N}$ 有 $\mu(A_n)<\infty$。
\item Borel 测度 $\mu$ 的支集（记为 $\operatorname{supp}\mu$）是闭集 $\{x\in X\mid\text{对 }x\text{ 的任意邻域 }U\text{ 有 }\mu(U)>0\}$。
\end{enumerate}
$X$ 上 Borel 测度与 Radon 测度的集合分别记为 $\mathcal{B}^+(X)$ 与 $\mathcal{R}^+(X)$。
\end{definition}

\begin{remark}\label{remark:ch01-Borel测度性质}
设 $X$、$\mathcal{A}$、$\mathcal{B}$、$\mu$ 如\cref{definition:ch01-Borel测度}所述。
\begin{enumerate}[(i)]
\item 集合 $E\in\mathcal{A}$ 是 $\mu$--外正则的，当且仅当对任意 $\epsilon>0$，存在 $A\in\mathcal{T}$ 使得 $A\supseteq E$ 且 $\mu(A\setminus E)<\epsilon$。集合 $E\in\mathcal{A}$ 是 $\mu$--内正则的，当且仅当对任意 $\epsilon>0$，存在 $K\in\mathcal{K}$ 使得 $K\subseteq E$ 且 $\mu(E\setminus K)<\epsilon$。若 $E\in\mathcal{A}$ 是 $\mu$--外正则的，则存在 $G\in\mathcal{B}$ 使得 $G\supseteq E$ 且 $\mu(G\setminus E)=0$。
\item $\sigma$--紧集 $E\subseteq X$ 是 $\mu$--内正则的。事实上，由定义存在非减序列 $\{K_n\}\subseteq\mathcal{K}$ 使得 $E=\bigcup_{n=1}^\infty K_n$，因而 $E\in\mathcal{B}$。于是由\cref{proposition:ch01-测度性质}(iv) 有
\[
\mu(E)=\lim_{n\to\infty}\mu(K_n)=\sup\{\mu(K)\mid K\subseteq E,\ K\in\mathcal{K}\}.
\]
\item 若 $\mu$ 局部有限（特别地，若 $\mu$ 为 Radon 测度），则对任意 $K\in\mathcal{K}$ 有 $\mu(K)<\infty$。事实上，对任意 $K\in\mathcal{K}$，存在 $A\in\mathcal{T}$ 使得 $A\supseteq K$ 且 $\mu(A)<\infty$，因而 $\mu(K)\leqslant\mu(A)<\infty$。
\item 若 Hausdorff 空间 $X$ 局部紧，则 $\mu$ 局部有限当且仅当对任意 $K\in\mathcal{K}$ 有 $\mu(K)<\infty$。事实上，"若"部分由局部紧空间的定义即得，而"仅当"部分由 (iii) 得到。特别地，若 $X$ 为局部紧 Hausdorff 空间，则 $\mu\in\mathcal{B}^+(X)$ 当且仅当对任意 $K\in\mathcal{K}$ 有 $\mu(K)<\infty$。
\item 每个适度 Borel 测度是 $\sigma$--有限的。反之，每个 $\sigma$--有限的外正则 Borel 测度是适度的。
\item 若 Hausdorff 空间 $X$ 为 $\sigma$--紧的，则 $X$ 上每个 Borel 测度 $\mu$ 都是适度的。事实上，设 $X=\bigcup_{n=1}^\infty K_n$，$\{K_n\}\subseteq\mathcal{K}$。由于 $\mu$ 局部有限，由 (iii) 对每个 $n\in\mathbb{N}$ 存在 $A_n\in\mathcal{T}$ 使得 $A_n\supseteq K_n$ 且 $\mu(A_n)<\infty$，故结论成立。
\end{enumerate}
\end{remark}

\begin{lemma}\label{lemma:ch01-有限Radon正则}\label{proposition:ch01-正则}
\label{lemma:ch01-紧外正则}
\label{lemma:ch01-正则性}
设 $(X,\mathcal{T})$ 为一 Hausdorff 空间，$\mu\in\mathcal{R}^+(X)$。则：
\begin{enumerate}[(i)]
\item 若 $\mu(X)<\infty$，则 $\mu$ 是正则的；
\item 每个 $K\in\mathcal{K}$ 是 $\mu$--外正则的。
\end{enumerate}
\end{lemma}

\begin{proof}
(i) 设 $E\in\mathcal{A}\supseteq\mathcal{B}$。由于 $\mu$ 内正则，集合 $E^c\in\mathcal{A}$ 是 $\mu$--内正则的，
\[
\mu(E^c)=\sup\{\mu(K)\mid K\subseteq E^c,\ K\in\mathcal{K}\}.
\]
因此
\[
\mu(E)=\mu(X)-\sup\{\mu(K)\mid K\subseteq E^c,\ K\in\mathcal{K}\}
=\inf\{\mu(X)-\mu(K)\mid K\subseteq E^c,\ K\in\mathcal{K}\}
=\inf\{\mu(K^c)\mid K^c\supseteq E,\ K\in\mathcal{K}\}.
\]
由于
\[
\mu(E)\leqslant\inf\{\mu(A)\mid A\supseteq E,\ A\in\mathcal{T}\}\leqslant\inf\{\mu(K^c)\mid K^c\supseteq E,\ K\in\mathcal{K}\},
\]
结论成立。

(ii) 固定 $K\in\mathcal{K}$。由于 $\mu$ 局部有限，如\cref{remark:ch01-Borel测度性质}(iii) 中所述，存在 $A\in\mathcal{T}$ 使得 $A\supseteq K$ 且 $\mu(A)<\infty$。由于 $\mu$ 内正则，集合 $A\setminus K\in\mathcal{B}$ 是 $\mu$--内正则的。于是对任意 $\epsilon>0$，存在 $L\in\mathcal{K}$、$L\subseteq A\setminus K$，使得
\[
\mu(L)>\mu(A\setminus K)-\epsilon.
\]
故 $M\triangleq A\setminus L\in\mathcal{T}$，$M\supseteq K$，且
\[
\mu(K)\leqslant\mu(M)=\mu(A)-\mu(L)<\mu(A)-\mu(A\setminus K)+\epsilon=\mu(K)+\epsilon.
\]
由 $\epsilon$ 的任意性，结论成立。
\end{proof}

\begin{proposition}\label{proposition:ch01-可数基Borel正则}
\label{remark:ch01-附录A1}
设 $(X,\mathcal{T})$ 为具有可数基的局部紧 Hausdorff 空间，$\mu\in\mathcal{B}^+(X)$。则 $\mu$ 是适度的且正则的。
\end{proposition}

\begin{proof}
由\cref{remark:ch01-附录A1}，$X$ 是 $\sigma$--紧的，故由\cref{remark:ch01-Borel测度性质}(vi)，$\mu$ 是适度的。

为证 $\mu$ 外正则，设 $\{A_n\}\subseteq\mathcal{T}$，对所有 $n\in\mathbb{N}$ 有 $\mu(A_n)<\infty$，且 $X=\bigcup_{n=1}^\infty A_n$。对每个 $E\in\mathcal{A}\supseteq\mathcal{B}$，令 $E_n\triangleq E\cap A_n$，于是 $E=\bigcup_{n=1}^\infty E_n$ 且对所有 $n\in\mathbb{N}$ 有 $\mu(E_n)\leqslant\mu(A_n)<\infty$。则 $\mu_n\triangleq \mu|_{\mathcal{B}(A_n)}$ 是有限的；此外，由\cref{remark:ch01-附录A1}，$A_n$ 的每个开子集是 $\sigma$--紧的，因而是 $\mu_n$--内正则的（见\cref{remark:ch01-Borel测度性质}(ii)）。如\cref{lemma:ch01-有限Radon正则}(i) 的证明那样论证，可知 $\mu_n$ 是正则的，特别地是外正则的（细节见文献~[45, Satz VIII.1.5]）。因此，对任意 $\epsilon>0$ 和 $n\in\mathbb{N}$，存在 $A_n'\in\mathcal{T}$ 使得
\[
E_n\subseteq A_n'\subseteq A_n,\qquad\mu_n(A_n'\setminus E_n)=\mu(A_n'\setminus E_n)<\frac{\epsilon}{2^n}.
\]
于是 $A'\triangleq \bigcup_{n=1}^\infty A_n'\in\mathcal{T}$，$A'\supseteq E$，且
\[
\mu(A'\setminus E)\leqslant\sum_{n=1}^\infty\mu(A_n'\setminus E_n)<\epsilon.
\]
故 $\mu$ 是外正则的（见\cref{remark:ch01-Borel测度性质}(i)）。

现证 $\mu$ 内正则。设 $E\in\mathcal{A}$，$E_n\triangleq E\cap A_n$，$\mu_n$ 同上。对任意 $\alpha<\mu(E)$，取 $p\in\mathbb{N}$ 充分大使得 $\mu(\bigcup_{n=1}^p E_n)>\alpha$；令 $\epsilon\triangleq \mu(\bigcup_{n=1}^p E_n)-\alpha$。对每个 $n=1,\cdots ,p$，测度 $\mu_n$ 是正则的，因而是内正则的；故存在紧集 $K_n\subseteq E_n$ 使得
\[
\mu_n(E_n\setminus K_n)=\mu(E_n\setminus K_n)<\frac{\epsilon}{p}\qquad(n=1,\cdots ,p).
\]
则 $K\triangleq \bigcup_{n=1}^p K_n\in\mathcal{K}$，$K\subseteq\bigcup_{n=1}^p E_n$，且
\[
\mu\left(\bigcup_{n=1}^p E_n\setminus K\right)\leqslant\sum_{n=1}^p\mu(E_n\setminus K_n)<\epsilon=\mu\left(\bigcup_{n=1}^p E_n\right)-\alpha.
\]
由上述不等式得 $\mu(K)>\alpha$，故结论成立。
\end{proof}

\begin{remark}\label{remark:ch01-BorelRadon一致}
由\cref{proposition:ch01-可数基Borel正则}，在具有可数基的局部紧 Hausdorff 空间中（例如在 $\mathbb{R}^N$ 中），Borel 测度与 Radon 测度重合且都是正则的。
\end{remark}

\subsection{零测集}\label{subsection:ch01-零测集}

\begin{definition}\label{definition:ch01-零测集}
设 $(X,\mathcal{A},\mu)$ 为一测度空间。
\begin{enumerate}[(i)]
\item 集合 $N\subseteq X$ 称为 $\mu$--零测的（或零测的），若 $N\in\mathcal{A}$ 且 $\mu(N)=0$。集合 $E\subseteq X$ 称为 $\mu$--可忽略的（或可忽略的），若存在 $\mu$--零测集 $N$ 使得 $E\subseteq N$。
\item 空间 $(X,\mathcal{A},\mu)$ 称为完备的，若所有可忽略集都是可测的，因而是零测的。此时 $\mu$ 是完备测度，$\sigma$--代数 $\mathcal{A}$ 关于 $\mu$ 是完备的。
\end{enumerate}
\end{definition}

我们记零测集族为 $\mathcal{N}_\mu$，可忽略集族为 $\mathcal{U}_\mu$。两个族 $\mathcal{N}_\mu$ 与 $\mathcal{U}_\mu$ 关于可数并都是稳定的。此外，$\mathcal{U}_\mu$ 是遗传的，即 $F\subseteq E$、$E\in\mathcal{U}_\mu\Rightarrow F\in\mathcal{U}_\mu$。显然，$\mathcal{N}_\mu=\mathcal{U}_\mu$ 当且仅当 $(X,\mathcal{A},\mu)$ 是完备的。

\begin{theorem}\label{theorem:ch01-Lebesgue完备化}\label{definition:ch01-完备化}
设 $(X,\mathcal{A},\mu)$ 为一测度空间。则：
\begin{enumerate}[(i)]
\item 族
\begin{align}
\overline{\mathcal{A}}\triangleq \{E\subseteq X\mid\exists F,G\in\mathcal{A}\text{ 使得 }F\subseteq E\subseteq G,\ \mu(G\setminus F)=0\} \label{align:ch01-完备sigma代数}
\end{align}
是包含 $\mathcal{A}$ 的 $\sigma$--代数；
\item 存在完备测度 $\overline{\mu}:\overline{\mathcal{A}}\to[0,\infty]$ 使得 $\overline{\mu}|_{\mathcal{A}}=\mu$。此外，$(X,\overline{\mathcal{A}},\overline{\mu})$ 是包含 $(X,\mathcal{A},\mu)$ 的最小完备测度空间。
\end{enumerate}
\end{theorem}

\begin{definition}\label{definition:ch01-Lebesgue完备化}
空间 $(X,\overline{\mathcal{A}},\overline{\mu})$ 称为 $(X,\mathcal{A},\mu)$ 的 Lebesgue 完备化。
\end{definition}

\begin{proof}[\cref{theorem:ch01-Lebesgue完备化}的证明]
(i) 显然 $\overline{\mathcal{A}}\supseteq\mathcal{A}$；特别地，$\varnothing \in\overline{\mathcal{A}}$。由\eqref{align:ch01-完备sigma代数} 立即得 $E\in\overline{\mathcal{A}}\Rightarrow E^c\in\overline{\mathcal{A}}$。此外设 $\{E_n\}\subseteq\overline{\mathcal{A}}$。则存在 $\{F_n\}\subseteq\mathcal{A}$ 和 $\{G_n\}\subseteq\mathcal{A}$ 使得对任意 $n\in\mathbb{N}$ 有 $F_n\subseteq E_n\subseteq G_n$，$\mu(G_n\setminus F_n)=0$。故
\[
F\triangleq \bigcup_{n=1}^\infty F_n\subseteq\bigcup_{n=1}^\infty E_n\subseteq\bigcup_{n=1}^\infty G_n=:G,
\]
由 $\mu$ 的 $\sigma$--次可加性有
\[
\mu(G\setminus F)\leqslant\sum_{n=1}^\infty\mu(G_n\setminus F_n)=0.
\]
故 $\bigcup_{n=1}^\infty E_n\in\overline{\mathcal{A}}$，结论成立。

(ii) 对任意 $E\in\overline{\mathcal{A}}$，令 $\overline{\mu}(E)\triangleq \mu(F)=\mu(G)$，其中 $F,G$ 如\eqref{align:ch01-完备sigma代数} 所示。注意此定义是适定的，因为它不依赖于 $F,G$ 的选取。事实上，设 $F_1,G_1\in\mathcal{A}$ 使得 $F_1\subseteq E\subseteq G_1$ 且 $\mu(G_1\setminus F_1)=0$。则 $F_1\subseteq E\subseteq G$ 且 $F\subseteq E\subseteq G_1$，故
\[
G\setminus F_1=(G\setminus E)\cup(E\setminus F_1)\subseteq(G\setminus F)\cup(G_1\setminus F_1).
\]
由此得 $\mu(G\setminus F_1)=0$；类似地有 $\mu(G_1\setminus F)=0$。等式
\[
\mu(G\setminus F)=\mu(G_1\setminus F_1)=\mu(G\setminus F_1)=\mu(G_1\setminus F)=0
\]
蕴含 $\mu(F)=\mu(F_1)=\mu(G_1)=\mu(G)$。

现证 $\overline{\mu}$ 是 $\overline{\mathcal{A}}$ 上的测度。显然 $\overline{\mu}(\varnothing )=0$。为证 $\overline{\mu}$ 的 $\sigma$--可加性，设 $\{E_n\}\subseteq\overline{\mathcal{A}}$ 为不交序列。于是每个满足 $F_n\subseteq E_n$（对所有 $n\in\mathbb{N}$）的序列 $\{F_n\}\subseteq\mathcal{A}$ 也是不交的。故由 $\mu$ 的 $\sigma$--可加性，
\[
\overline{\mu}\left(\bigcup_{n=1}^\infty E_n\right)=\mu\left(\bigcup_{n=1}^\infty F_n\right)=\sum_{n=1}^\infty\mu(F_n)=\sum_{n=1}^\infty\overline{\mu}(E_n),
\]
结论成立。为证 $\overline{\mu}$ 完备，固定任意 $N\in\overline{\mathcal{A}}$ 使得 $\overline{\mu}(N)=0$，并设 $F,G\in\mathcal{A}$ 满足 $F\subseteq N\subseteq G$，$\mu(G\setminus F)=0$。由 $\overline{\mu}$ 的定义有 $\mu(G)=\overline{\mu}(N)=0$。于是对任意 $E\subseteq N$，有 $\varnothing \subseteq E\subseteq G$，$\mu(G\setminus\varnothing )=\mu(G)=0$。这表明 $E\in\overline{\mathcal{A}}$，故 $\overline{\mu}$ 完备。

由其定义有 $\overline{\mu}|_{\mathcal{A}}=\mu$。还需证：若 $(X,\mathcal{A}_1,\mu_1)$ 为完备测度空间且 $\mathcal{A}\subseteq\mathcal{A}_1$、$\mu_1|_{\mathcal{A}}=\mu$，则 $\overline{\mathcal{A}}\subseteq\mathcal{A}_1$ 且 $\mu_1|_{\overline{\mathcal{A}}}=\overline{\mu}$。为此，设 $E\in\overline{\mathcal{A}}$，$F,G\in\mathcal{A}$ 满足 $F\subseteq E\subseteq G$ 且 $\mu(G\setminus F)=0$。由于 $\mu_1$ 完备，$\mu_1(G\setminus F)=\mu(G\setminus F)=0$，且 $G\setminus E\subseteq G\setminus F$，故 $G\setminus E\in\mathcal{A}_1$。于是 $E=G\setminus(G\setminus E)\in\mathcal{A}_1$，因而 $\overline{\mathcal{A}}\subseteq\mathcal{A}_1$。其余结论显然，证毕。
\end{proof}

\begin{remark}\label{remark:ch01-完备化刻画}
值得注意\eqref{align:ch01-完备sigma代数} 中定义的 $\sigma$--代数 $\overline{\mathcal{A}}$ 的如下刻画：
\[
\overline{\mathcal{A}}=\{F\cup E_0\mid F\in\mathcal{A},\ E_0\in\mathcal{U}_\mu\}.
\]
事实上，设 $E\in\overline{\mathcal{A}}$。则由\eqref{align:ch01-完备sigma代数} 有 $E_0\triangleq E\setminus F\subseteq G\setminus F$，其中 $G\setminus F\in\mathcal{N}_\mu$；故 $E_0\in\mathcal{U}_\mu$ 且 $E=F\cup E_0$。反之，设 $E=F\cup E_0$，$F\in\mathcal{A}$，$E_0\in\mathcal{U}_\mu$。则由 $\mathcal{U}_\mu$ 的定义存在 $N_0\in\mathcal{N}_\mu$ 使得 $E_0\subseteq N_0$。令 $G\triangleq F\cup(N_0\setminus F)$。则有 $F\subseteq E\subseteq G$，$\mu(G\setminus F)=\mu(N_0\setminus F)=0$，故 $E\in\overline{\mathcal{A}}$。
\end{remark}

\begin{definition}\label{definition:ch01-a.e.}
设 $(X,\mathcal{A},\mu)$ 为一测度空间。性质 $P$ 称为 $\mu$--几乎处处成立（记为 $\mu$--a.e. 或 a.e.），若集合 $\{x\in X\mid P(x)\text{ 不成立}\}\in\mathcal{N}_\mu$。
\end{definition}

作为上述定义的特殊情形，(a) 两个函数 $f,g:X\to\mathbb{R}$ 在 $X$ 上 $\mu$--a.e. 相等，若 $\{x\in X\mid f(x)\neq g(x)\}\in\mathcal{N}_\mu$；(b) 序列 $\{f_n\}$（$f_n:X\to\mathbb{R}$）在 $X$ 中 $\mu$--a.e. 收敛，若 $\{x\in X\mid\{f_n(x)\}\text{ 不收敛}\}\in\mathcal{N}_\mu$。

显然，$\mu$--a.e. 相等是一个等价关系：
\begin{align}
g\sim f \overset{\mathrm{def}}{\Longleftrightarrow} g=f\ \mu\text{--a.e. 在 }X\text{ 中}. \label{align:ch01-a.e.等价}
\end{align}
关于\eqref{align:ch01-a.e.等价} 的每个等价类由其任意一个元素唯一确定，该元素称为该类的一个代表元。因此自然地将整个类视为一个在 $X$ 中 $\mu$--a.e. 定义的唯一映射 $f$。

\begin{definition}\label{definition:ch01-a.e.定义}
设 $(X,\mathcal{A},\mu)$ 为一测度空间，$Y$ 为一集合。函数 $f:E\to Y$（$E\in\mathcal{A}$）称为在 $X$ 中 $\mu$--a.e. 定义的，若 $E^c\in\mathcal{N}_\mu$。
\end{definition}

在适当场合，为强调\cref{definition:ch01-a.e.定义} 与经典逐点定义之间的区别，我们说：
\begin{enumerate}[(i)]
\item 函数 $f:X\to\mathbb{R}$ 是经典的，若它在每个点 $x\in X$ 处都有定义；
\item 经典函数 $g$ 等于一个 $\mu$--a.e. 定义的函数 $f$，若 $g$ 是 $f$ 的一个经典代表元。
\end{enumerate}

在此联系下，值得注意序列的 $\mu$--a.e. 极限在\cref{definition:ch01-a.e.定义} 的意义下是唯一的（但在逐点意义下不唯一）：事实上，$f_n\to f$ $\mu$--a.e. 且 $f_n\to g$ $\mu$--a.e. 蕴含 $f=g$ $\mu$--a.e.。

为后续参考，我们证明以下技术引理。

\begin{lemma}[穷尽引理]\label{lemma:ch01-穷尽引理}
\label{lemma:ch01-穷竭}
设 $(X,\mathcal{A},\mu)$ 为一有限测度空间。设 $P=P(x)$（$x\in X$）为一性质，满足：
\begin{enumerate}[(i)]
\item $P$ 在 $X$ 的每个 $\mu$--零测子集上成立；
\item 对任意 $E\in\mathcal{A}$ 且 $\mu(E)>0$，存在 $F\in\mathcal{A}$、$F\subseteq E$、$\mu(F)>0$，使得 $P$ 在 $F$ 上成立。
\end{enumerate}
则存在不交序列 $\{E_k\}\subseteq\mathcal{A}$ 使得 $\bigcup_{k=1}^\infty E_k=X$ 且 $P$ 在每个 $E_k$ 上成立。
\end{lemma}

\begin{proof}
令 $\Sigma\triangleq \{E\in\mathcal{A}\mid P\text{ 在 }E\text{ 上成立}\}$，$c\triangleq \sup_{E\in\Sigma}\mu(E)$。设 $\{F_n\}\subseteq\Sigma$ 使得 $\lim_{n\to\infty}\mu(F_n)=c$。令 $H_k\triangleq \bigcup_{n=1}^k F_n$。则对任意 $k\in\mathbb{N}$，性质 $P$ 在 $H_k$ 上成立，且 $H_k\subseteq H_{k+1}$；此外有 $\lim_{k\to\infty}\mu(H_k)=c$。

现证集合 $E_1\triangleq X\setminus(\bigcup_{k=1}^\infty H_k)=\bigcap_{k=1}^\infty H_k^c$ 是 $\mu$--零测的。反设 $\mu(E_1)>0$。则由假设 (ii) 存在 $F\in\mathcal{A}$、$F\subseteq E_1$、$\mu(F)>0$，使得性质 $P$ 在 $F$ 上成立，故 $F\in\Sigma$；此外对所有 $k\in\mathbb{N}$ 有 $F\subseteq H_k^c$。令 $G_k\triangleq H_k\cup F$（$k\in\mathbb{N}$）。则 $\{G_k\}\subseteq\Sigma$ 且
\[
\lim_{k\to\infty}\mu(G_k)=\lim_{k\to\infty}(\mu(H_k)+\mu(F))=c+\mu(F)>c,
\]
矛盾。

现令 $E_2\triangleq H_1$，$E_k\triangleq H_{k-1}\setminus H_{k-2}$ 对任意 $k\geqslant 3$。序列 $\{E_k\}$ 具有所述性质，证毕。
\end{proof}

最后我们证明以下关于 Radon 测度支集的刻画。

\begin{lemma}\label{lemma:ch01-支集刻画}
设 $(X,\mathcal{T})$ 为一 Hausdorff 空间，$\mu\in\mathcal{R}^+(X)$。设 $A\in\mathcal{T}$ 为最大的开 $\mu$--零测集。则 $\operatorname{supp}\mu=A^c$。
\end{lemma}

\begin{proof}
我们首先证明最大的开 $\mu$--零测集确实存在。令 $\mathcal{N}_\mu\cap\mathcal{T}\equiv\{A_j\}_{j\in J}$，$A\triangleq \bigcup_{j\in J}A_j$。显然 $A\in\mathcal{T}$ 且对所有 $j\in J$ 有 $A\supseteq A_j$。若 $\mathcal{N}_\mu\cap\mathcal{T}=\{\varnothing \}$，则结论显然。否则，设 $K\subseteq A$ 为紧集。于是存在 $A_1,\cdots ,A_n\in\mathcal{N}_\mu\cap\mathcal{T}$ 使得 $K\subseteq\bigcup_{j=1}^n A_j$。由此得 $\mu(K)=0$，由 $K$ 的任意性和 $\mu$ 的内正则性得 $\mu(A)=0$。

设 $x\in[\operatorname{supp}\mu]^c$。则由\cref{definition:ch01-Borel测度}(g) 存在开邻域 $\widetilde{A}\ni x$ 使得 $\mu(\widetilde{A})=0$；特别地，$\widetilde{A}\in\mathcal{N}_\mu\cap\mathcal{T}$。故 $x\in\widetilde{A}\subseteq A$。反之，设 $x\in A$。则 $A$ 是 $x$ 的一个开 $\mu$--零测邻域，故 $x\in[\operatorname{supp}\mu]^c$。于是 $[\operatorname{supp}\mu]^c=A$，证毕。
\end{proof}




\end{document}